\section{语言模型智能体：基本概念}

本讲的第二部分转向语言模型智能体（Language Model Agents）。智能体与推理密切相关：在环境中完成任务需要对后续状态、对象可操作性和不确定性进行\textbf{多步推理和规划}。

\subsection{智能体的核心术语}

\begin{figure}[H]
\centering
\includegraphics[width=0.95\textwidth]{slides-images/slide-040.jpg}
\caption{智能体的核心术语：Agent 从 Environment 接收 Observation，基于 Language Instruction ($g$) 发出 Action\protect\footnotemark}
\end{figure}
\footnotetext{Slides 第41页。}

\begin{importantbox}{智能体的三要素}
\begin{itemize}
    \item \textbf{Agent（智能体）}：执行任务的神经网络（或语言模型），其策略 $\pi(\cdot | g)$ 以语言指令 $g$ 为条件
    \item \textbf{Environment（环境）}：智能体交互的对象，可以是网页浏览器、数据库、物理世界的模拟等
    \item \textbf{三个核心变量}：Observation（环境观测，如 HTML DOM 或屏幕截图）、Action（动作，如点击、输入文字）、Language Instruction（自然语言指令，如``订一张从旧金山到纽约的机票''）
\end{itemize}
\end{importantbox}

这种设置有多种名称：数字智能体（Digital Agent）、语言条件策略（Language Conditioned Policy）、指令跟随智能体（Instruction-Following Agent）。

\subsection{应用场景}

语言模型智能体的应用场景极为广泛：

\begin{itemize}
    \item \textbf{数字助理}：自然语言控制手机设置闹钟、设置提醒等
    \item \textbf{自然语言编程}：根据自然语言描述生成 Python 代码
    \item \textbf{UI 自动化测试}：用语言模型替代人工测试 UI 元素
    \item \textbf{用户场景自动化}：如在 Spotify 上播放歌曲、在电商网站购物
    \item \textbf{工具使用}：为语言模型添加插件或工具，使其能控制各种应用
\end{itemize}

\begin{figure}[H]
\centering
\includegraphics[width=0.95\textwidth]{slides-images/slide-042.jpg}
\caption{智能体的应用场景：包括数字助理、自然语言编程、UI自动化和工具使用等\protect\footnotemark}
\end{figure}
\footnotetext{Slides 第43页。}

\subsection{本章小结}

语言模型智能体将 LLM 从纯文本生成扩展到环境交互，核心框架是：Agent 根据 Observation 和 Language Instruction 生成 Action，通过与 Environment 的循环交互来完成任务。这种框架统一了从网页操作到代码生成的各种应用场景。

%% ========== Part 6: LLM 之前的方法 ==========

